\documentclass[11pt]{amsart}
\usepackage[T1]{fontenc}
\usepackage{lmodern}
\usepackage{amsmath,amssymb}
\usepackage{microtype}
\usepackage[hidelinks,pdftitle={Draft: Paired clique minors from connected regions},pdfauthor={Cavit Erginsoy}]{hyperref}
\setlength{\emergencystretch}{1em}
\raggedbottom

\newtheorem{theorem}{Theorem}[section]
\newtheorem{lemma}[theorem]{Lemma}
\newtheorem{corollary}[theorem]{Corollary}
\newtheorem{proposition}[theorem]{Proposition}

\title[Paired clique minors]{Paired clique minors from connected regions}
\author{Cavit Erginsoy}
\date{Draft, 21 September 2026}

\begin{document}
\begin{abstract}
Let $R,S$ be disjoint sets of $k$ vertices joined by $k$ vertex-disjoint
paths. We prove that $k-1$ disjoint connected regions, each adjacent to
every vertex of $R\cup S$ and avoiding the terminals, force a $K_k$-minor
whose branch sets each contain one vertex of $R$ and one of $S$.
The pairing is not prescribed, and the number of regions is best
possible. In a stronger one-sided theorem, the regions cover $S$ and
each meets $S$, while adjacency is required only to $R$. A chain of
minimum vertex separators shows that a linkage in a reduced graph has
pairwise adjacent paths. The reductions give a polynomial-time
construction when the regions are supplied.
\end{abstract}
\maketitle

\section{Statements}

All graphs are finite and simple. For disjoint vertex sets $R,S$,
an \emph{$R$--$S$ linkage of order $k$} is a family of $k$ pairwise
vertex-disjoint paths with distinct ends in $R$ and in $S$. We write
$\kappa_G(R,S)$ for the largest possible order. A vertex set is
\emph{full to} $R$ if every vertex of $R$ has a neighbour in it.
A \emph{region} is a nonempty vertex set inducing a connected subgraph.

For disjoint $k$-element sets $R,S$, a \emph{paired clique model} consists
of $k$ pairwise disjoint nonempty connected branch sets that are pairwise
adjacent, each containing exactly one vertex of $R$ and exactly one
vertex of $S$. This is a $K_k$-minor model with a free terminal pairing.
The regions in the following statements need not be preserved as
additional branch sets.

\begin{theorem}\label{thm:one-sided}
Let $k\geq2$, let $R,S$ be disjoint $k$-element vertex sets in a graph
$G$, and suppose that $\kappa_G(R,S)\geq k$. Suppose there are $k-1$
pairwise disjoint regions $H_1,\ldots,H_{k-1}$ avoiding $R$, each full
to $R$, such that
\[
 S\subseteq\bigcup_{j=1}^{k-1}H_j
 \quad\text{and}\quad S\cap H_j\ne\varnothing\quad(1\leq j<k).
\]
Then $G$ has a paired clique model for $R,S$.
\end{theorem}

The theorem requires no global connectivity assumption on $G$ and implies
the following symmetric statement.

\begin{corollary}\label{cor:two-sided}
Let $k\geq1$ and let $R,S$ be disjoint $k$-element vertex sets in $G$.
Suppose $G$ contains $k-1$ pairwise disjoint regions avoiding $R\cup S$,
each full to $R\cup S$. Then $G$ has a paired clique model for $R,S$
if and only if $\kappa_G(R,S)\geq k$.
\end{corollary}

We prove Theorem~\ref{thm:one-sided} in
Sections~\ref{sec:reductions} through~\ref{sec:contacts}.
Section~\ref{sec:consequences} gives the polynomial-time construction
and shows that the number of regions in Corollary~\ref{cor:two-sided}
cannot be reduced below $k-1$, for any $k\geq2$.

In the terminology of Protopapas, Thilikos and
Wiederrecht~\cite[Sections 1.1 and 2.1]{PTW}, the conclusion is a
\emph{rainbow $K_k$-minor} when $R$ and $S$ are the two colour sets
and all other vertices are uncoloured. Each branch set meets both
colour sets exactly once, since each set has $k$ vertices.
Their Theorem~3.1, specialised to two colours and a rainbow
$K_k$, starts from a $K_{4k}$-minor and gives a rainbow minor or a
separator outcome. Our sufficient condition uses the supplied connected
regions and the linkage. We make no claim of publication priority.

An \emph{$R$--$S$ separator} is a set $C\subseteq V(G)$ such that
$G-C$ has no path from $R\setminus C$ to $S\setminus C$; terminal
vertices are allowed in $C$. We use the vertex form of Menger's
theorem~\cite{Menger}: the minimum size of such a separator equals
$\kappa_G(R,S)$. In particular, when $|R|=|S|=k$ and a linkage has
order $k$, every terminal is a path end and no terminal is internal
to a linkage path.

\section{Reductions and terminal ownership}\label{sec:reductions}

Suppose Theorem~\ref{thm:one-sided} is false. Choose a counterexample
with $k$ minimum and, subject to this, with $|V(G)|$ minimum. Since
$k$ vertices of $S$ occupy $k-1$ regions, each meeting $S$, exactly
one region contains two of them and every other region contains one.

No region is a singleton. Indeed, a singleton region is $\{s\}$ for
some $s\in S$. Fix a linkage of order $k$ and let $r\in R$ be the
other end of its path ending at $s$. Delete $r,s$. The other $k-1$
paths form a linkage for $R\setminus\{r\},S\setminus\{s\}$, and the
remaining $k-2$ regions satisfy the hypotheses of the theorem for
$k-1$. Here $k\geq3$: when $k=2$, the only region contains both
vertices of $S$. Minimality of $k$ supplies a paired clique model in
$G-\{r,s\}$. Add the branch set $\{r,s\}$. It is connected, since
$\{s\}$ is full to $R$, and $s$ is adjacent to the $R$ vertex in each
other branch set. The earlier branch sets avoid $r,s$, since both
terminals were deleted before induction.

We may assume that the regions cover $V(G)\setminus R$. To see this,
consider a component $D$ of
$G-(R\cup\bigcup_j H_j)$. If $D$ contacts a region, absorb it into
that region. This preserves disjointness, connectivity, fullness and
all terminal memberships. If it contacts no region, its neighbourhood
is contained in $R$. Every linkage of order $k$ avoids $D$: a path
visiting $D$ would have to pass through an $R$ terminal internally to
reach its $S$ end. Deleting $D$ therefore preserves the hypotheses,
contrary to minimum host order. Absorbing all remaining components gives
\begin{equation}\label{eq:partition}
 V(G)=R\mathbin{\dot\cup}H_1\mathbin{\dot\cup}\cdots
       \mathbin{\dot\cup}H_{k-1}.
\end{equation}
The graph has not changed, and no region is a singleton.

Call an edge within a region \emph{permitted} if its ends are not both
in $S$. Contracting a permitted edge preserves the two terminal sets:
an $S$ endpoint, if present, retains its name. The resulting regions
are still disjoint and connected, full to $R$, and have the same $S$
memberships. If the quotient has a linkage of order $k$, minimum host
order supplies its paired clique model. Lift that model by replacing
the contracted vertex, if used, by its connected two-vertex preimage.
This preimage contains no $R$ vertex and at most one $S$ vertex. The
lift preserves disjointness, all contacts and every terminal assignment.
Consequently,
\begin{equation}\label{eq:failed-contraction}
 \kappa_{G/uv}(R,S)\leq k-1
 \quad\text{for every permitted region edge }uv.
\end{equation}

\section{Separators and a spanning linkage}\label{sec:separators}

\begin{lemma}\label{lem:edge-cut}
For every permitted edge $uv$ in $H_j$, there is an $R$--$S$ separator
of size $k$ containing $u,v$, exactly one vertex in each other region,
and no vertex of $R$. For any fixed linkage of order $k$, these $k$
separator vertices lie one on each path.
\end{lemma}

\begin{proof}
By~\eqref{eq:failed-contraction} and Menger's theorem, $G/uv$ has a
separator $C'$ of size at most $k-1$. It contains the contracted
vertex, since otherwise $C'$ would also separate $R$ from $S$ in $G$.
Replace that vertex by $u,v$ to obtain a separator $C$ in $G$ of size
at most $k$. The existing linkage forces $|C|=k$.

Since $u,v\notin R$, the cut $C$ removes at most $k-2$ vertices of
$R$, so some $R$ vertex survives. If $C$ missed another region, that
region would connect this vertex to one of its $S$ terminals in $G-C$.
The remaining $k-2$ cut vertices must therefore lie one in each of
the other $k-2$ regions. In particular, $C$ avoids $R$. Each of the
$k$ disjoint linkage paths meets $C$, so $|C|=k$ gives exactly one
meeting per path.
\end{proof}

Fix a linkage $P_1,\ldots,P_k$ and orient each path from its $R$ end
to its $S$ end. This linkage spans $G$: every region vertex outside $S$
has a neighbour in its region because the region is connected and
non-singleton. The incident edge is permitted, and
Lemma~\ref{lem:edge-cut} places that vertex on the linkage. All
terminals are already path ends, so~\eqref{eq:partition} proves the
claim. The ends of every region edge lie on different paths:
use Lemma~\ref{lem:edge-cut} for a permitted edge, and the distinct
$S$ ends for an edge with both ends in $S$.

\section{A chain of separators}\label{sec:chain}

Write
\[
 P_i=p_{i,0}p_{i,1}\cdots p_{i,\ell_i},\qquad
 p_{i,0}\in R,\quad p_{i,\ell_i}\in S.
\]
Consider all vectors $c=(c_1,\ldots,c_k)$ with
$1\leq c_i\leq\ell_i$ such that
$\{p_{i,c_i}:1\leq i\leq k\}$ is an $R$--$S$ separator.
Identify each vector with its cut and order the vectors coordinatewise.
The cuts exclude $R$ endpoints but may contain $S$ endpoints. The vector
$(\ell_1,\ldots,\ell_k)$ selects $S$ and is the maximum of the collection.

Since the linkage spans $G$, a vector $c$ is separating exactly when
there is no edge from a strict prefix
$\{p_{i,t}:t<c_i\}$ to a strict suffix
$\{p_{j,t}:t>c_j\}$. Every prefix vertex connects to its $R$ end
without meeting the selected cut, and every suffix vertex connects to
its $S$ end. A crossing edge would therefore give a path avoiding the
cut; without one, no surviving path can pass from a prefix to a suffix.
This remains true when a selected $S$ endpoint has an empty suffix.

The separating vectors are closed under coordinatewise minimum and
maximum, denoted by $\wedge$ and $\vee$. For a crossing edge at
$c\wedge d$, its prefix end is before both original cuts. Choose the
original cut attaining the minimum at the suffix end's coordinate;
the edge crosses that original cut, a contradiction. For $c\vee d$,
choose the original cut attaining the maximum at the prefix end's
coordinate. Its suffix end is after both original cuts, giving the
same contradiction.

Choose a maximal chain of separating vectors,
\[
 c^0<c^1<\cdots<c^m.
\]
It includes the minimum and maximum of this finite lattice. Every
region vertex occurs in a cut on the chain. For a vertex outside $S$,
Lemma~\ref{lem:edge-cut} provides a cut containing it; for an $S$
vertex, use the maximum cut. If a vertex $z=p_{i,t}$ were skipped,
the extreme cuts imply that two consecutive chain cuts $c<d$ would
satisfy $c_i<t<d_i$. A cut $a$ containing $z$ would then give the
strictly intermediate separating vector
\[
 (a\vee c)\wedge d,
\]
contrary to maximality of the chain. Monotonicity shows that the
indices of chain cuts containing $z$ form a nonempty interval $I_z$
of integers.

\begin{lemma}\label{lem:edge-overlap}
For every edge $uv$ within a region, $I_u\cap I_v\ne\varnothing$.
\end{lemma}

\begin{proof}
If $u,v\in S$, their intervals contain the maximum cut. Otherwise
Lemma~\ref{lem:edge-cut} gives a cut $a$ containing both ends, which
lie on different paths. Suppose the occurrence intervals are disjoint.
After interchanging the ends, let $r$ be the last index in $I_u$ and
$l$ the first in $I_v$, with $r<l$. At $c^{r+1}$, the vertex $u$ is
in a strict prefix. The actual edge $uv$ prevents $v$ from being in
a strict suffix, so $l=r+1$: this next cut contains $v$.

The separating vector
\[
 (a\vee c^r)\wedge c^{r+1}
\]
lies between the two consecutive cuts and contains both $u$ and $v$.
It differs from $c^r$ at $v$'s coordinate and from $c^{r+1}$ at
$u$'s coordinate, contradicting maximality. This step uses the edge
$uv$ to exclude $v$ from a strict suffix when $u$ enters a strict prefix.
\end{proof}

\section{The first cut gives all clique contacts}\label{sec:contacts}

Every chain cut meets every region. Indeed, the cut avoids $R$; an
untouched region would contain an uncut $S$ vertex and connect it to
$R$ by fullness. A cut has $k$ vertices and there are $k-1$ regions,
so each region contributes at most two vertices to a chain cut.

Fix a region $H_j$ and form the intersection graph $T_j$ of the
intervals $(I_z:z\in H_j)$. No three of these intervals contain a
common index. The graph $T_j$ is a forest: if it had a cycle, choose
on that cycle an interval with smallest right endpoint. Both its
cycle neighbours contain that endpoint, because each meets the chosen
interval and has a right endpoint no smaller. This would give three
intervals with a common index.

By Lemma~\ref{lem:edge-overlap}, the connected graph $G[H_j]$ is a
spanning subgraph of $T_j$. A forest with a connected spanning subgraph
is a tree and cannot have any further edge. Hence $G[H_j]=T_j$.
Thus two vertices of one region that occur in a common chain cut
are adjacent in $G$.

For each $i$, put $f_i=p_{i,1}$, the first vertex after the $R$ end
of $P_i$; it may be an $S$ end. Each $f_i$ occurs in a separating
vector. The coordinatewise minimum of one such vector for each $i$
is therefore
\[
 F=\{f_1,\ldots,f_k\}.
\]
This is the minimum separating vector, so it is on the chain.
Every region neighbour of an $R$ terminal lies in $F$. A neighbour
later on its path would have a suffix to $S$ avoiding $F$, giving
an $R$--$S$ path avoiding this cut.

Exactly one region is represented twice in $F$, and every other
region once. If $f_i,f_j$ lie in the same region, the preceding
tree argument gives the edge $f_if_j$. If they lie in different
regions, at least one of those regions is represented only once in
$F$. The $R$ endpoint of the other path is full to that region,
and all its neighbours there lie in $F$, so it is adjacent to that
unique representative. Thus $P_i,P_j$ are adjacent in every case.

The vertex sets of $P_1,\ldots,P_k$ form the required paired clique
model, a contradiction. This proves Theorem~\ref{thm:one-sided}.
The singleton step decreases $k$ after deleting both reserved terminals.
Contraction decreases host order and lifts through a fixed connected
preimage containing at most one terminal. Both steps preserve terminal
ownership. \qed

\section{The symmetric statement, sharpness and construction}
\label{sec:consequences}

\begin{proof}[Proof of Corollary~\ref{cor:two-sided}]
A paired clique model supplies a linkage by choosing an $R$--$S$
path inside each branch set. Conversely, suppose a linkage of order
$k$ exists. For $k=1$, take its only path as the branch set. For
$k\geq2$, assign the $k$ vertices of $S$ to the $k-1$ given regions,
at least one to each, and adjoin each vertex to its assigned region.
Fullness to $S$ makes the enlarged regions connected. They remain
pairwise disjoint, avoid $R$ and are full to $R$. They cover $S$ and
each meets $S$, so Theorem~\ref{thm:one-sided} applies.
\end{proof}

\begin{proposition}\label{prop:sharp}
For every $k\geq2$, Corollary~\ref{cor:two-sided} fails if its
number of full regions is reduced to $k-2$.
\end{proposition}

\begin{proof}
Take $k$ disjoint edges $r_is_i$ with
$R=\{r_1,\ldots,r_k\}$ and $S=\{s_1,\ldots,s_k\}$, and add
$k-2$ vertices, each adjacent to all $2k$ terminals. Add no other
edges. The matching is a linkage of order $k$, and each added vertex
is a full singleton region avoiding the terminals.

In any proposed paired clique model, at least two of its $k$ branch
sets contain no added vertex. Each of these consists of exactly one
$R$ vertex and one $S$ vertex. Connectivity forces it to be one of
the matching edges, and two distinct such branch sets are not
adjacent. This is a contradiction, including when $k=2$.
\end{proof}

\begin{corollary}\label{cor:algorithm}
Under the hypotheses of Theorem~\ref{thm:one-sided} or
Corollary~\ref{cor:two-sided}, a paired clique model can be found in
polynomial time when the regions are supplied.
\end{corollary}

\begin{proof}
For the symmetric statement, make the terminal assignment from the
proof of Corollary~\ref{cor:two-sided}; for $k=1$, return a path.
It remains to give a construction under the one-sided hypotheses.
Vertex-capacitated maximum flow finds a linkage of order $k$ and
tests whether one survives a specified contraction.

First absorb or delete components outside $R$ and the regions, as
in Section~\ref{sec:reductions}. If a region is a singleton
$\{s\}$, take a linkage of order $k$ and let $r$ be the other end
of its path ending at $s$. Recurse after deleting $r,s$ and this
region, then append the branch set $\{r,s\}$. This recursion has
$k-1\geq2$ and satisfies the smaller theorem's hypotheses.

Otherwise test every permitted region edge. If contracting one
preserves a linkage of order $k$, recurse on the quotient and lift
the returned branch sets through that edge. If all these tests fail,
return any linkage of order $k$. Its paths are pairwise adjacent by
Sections~\ref{sec:separators} through~\ref{sec:contacts}: after normalisation
and removal of the singleton case, those sections use minimality
only through~\eqref{eq:failed-contraction}. The algorithm therefore
does not construct the cut lattice or enumerate its maximal chains.

Every recursive call decreases the number of vertices. There are
therefore at most $|V(G)|$ calls, each using polynomially many flow
tests and elementary graph operations. The disjoint connected
preimages used for every contraction lift are recorded during the
recursion. Component deletions and the singleton step require no
further allocation, so all lifts also take polynomial time.
\end{proof}

\begin{thebibliography}{9}
\raggedright
\bibitem{Menger}
K.~Menger,
\href{https://doi.org/10.4064/fm-10-1-96-115}{\emph{Zur allgemeinen
Kurventheorie}}, Fund.~Math.\ \textbf{10} (1927), 96--115.

\bibitem{PTW}
E.~Protopapas, D.~M.~Thilikos and S.~Wiederrecht,
\emph{Colorful Minors},
\href{https://arxiv.org/abs/2507.10467v3}{arXiv:2507.10467v3}, 2026.
\end{thebibliography}
\end{document}
